% Generated draft: requires mathematical and source review.
\documentclass[11pt]{article}
\usepackage[margin=25mm]{geometry}
\usepackage{fontspec}
\setmainfont{Latin Modern Roman}
\usepackage{amsmath,amssymb,mathtools,unicode-math}
\setmathfont{Latin Modern Math}
\usepackage{enumitem,needspace,xurl,hyperref}
\setlength{\parindent}{0pt}
\setlength{\parskip}{5pt}
\setlength{\emergencystretch}{3em}
\newcommand{\N}{\mathbb{N}}
\newcommand{\Z}{\mathbb{Z}}
\newcommand{\Q}{\mathbb{Q}}
\newcommand{\R}{\mathbb{R}}
\newcommand{\C}{\mathbb{C}}
\newcommand{\F}{\mathbb{F}}
\newcommand{\A}{\mathbb{A}}
\newcommand{\PP}{\mathbb{P}}
\newcommand{\E}{\mathbb{E}}
\newcommand{\eps}{\varepsilon}
\newcommand{\dd}{\,\mathrm{d}}
\begin{document}
\section*{UnsolvedMath ID: 20002326}
\textbf{Source problem code:} catalogue entry
\subsection{Definitions and mathematical statement}
\textbf{Problem 1.} Let $T$ be a complete superstable metric theory, expand its language by an automorphism symbol $\sigma$, and put $T_\sigma=T\cup\{\sigma\text{ is an automorphism}\}$. Assume a model companion $T_A$ exists. Is $T_A$ supersimple up to perturbations of the automorphism? The source says ``unitary function''; for a general metric theory the intended extra symbol appears to be \emph{unary}, not a Hilbert-space unitary operator. The perturbation system and exact meaning of perturbative supersimplicity are not fixed. If the model companion is incomplete, the role of its completions must also be specified.

The source contrasts the classical superstable-to-supersimple result of Pillay--Chatzidakis with probability algebras: it reports that $T$ is superstable and $T_A$ exists but is not superstable (Ben-Yaacov), while it is superstable up to automorphism perturbations (Ben-Yaacov, Berenstein). These are supplied background examples, not a general transfer theorem.

\textbf{Supplied concrete perturbation estimate.} Let $(M,\sigma)$ and $(N,\tau)$ have an isomorphism $f:M\to N$ of their underlying $T$-structures, and suppose
\[
 \sup_{x\in M}d(f\sigma x,\tau f x)\le r.
\]
The automorphisms are surjective isometries. The aid asserts, for every integer $k$,
\[
 d(f\sigma^kx,\tau^kf x)\le|k|r.
\]
For an orbit formula $\varphi$ obtained from a base-language formula $\theta$ using orbit coordinates of powers $|k|\le N_0$, let $\omega_\theta$ be a uniform-continuity modulus in the maximum tuple metric. The corresponding transported formula values differ by at most
\[
 \omega_\theta(N_0r).
\]
The estimate is uniform in variable values and transported parameters; it assumes an exact underlying $T$-isomorphism, not an arbitrary approximately structure-preserving map.

Using continuous truth value zero, say a family $\{\varphi(x;a_i):i\in I\}$ has a uniform $m,\delta$ inconsistency gap when for every $m$ distinct indices,
\[
 \inf_x\max_{1\le j\le m}\varphi(x;a_{i_j})\ge\delta>0.
\]
Under the perturbation above the aid's finite-horizon bound gives the lower gap
\[
 \delta-\omega_\theta(N_0r).
\]
This is a positive surviving inconsistency gap only when $\omega_\theta(N_0r)<\delta$. Thus bounded-depth positive-gap patterns cannot be erased by arbitrarily small perturbations in this concrete system. Exact expanded-language indiscernibility and compatibility of all orbit data are not supplied by a finite-horizon estimate. The aid identifies these global issues as the remaining bottleneck and does not derive perturbative supersimplicity from the bound.

\textbf{Appended Problem 2.} The same extract independently asks (a) for metric analogues of Ehrenfeucht's non-countably-categorical theories with finitely many countable models, and (b) what follows when separable models, equipped with a chosen perturbation pseudometric, have a compact or finite quotient after identifying distance-zero models. A typical definition is $d_\rho(M,N)=\inf\{r:\text{an }r\text{-perturbation identifies }M,N\}$ for a fixed system $\rho$; the system, permitted models and equivalences must be specified for the question. Whether genuinely metric examples rather than discrete reinterpretations are intended is left open by the text. The terminal ``Background:'' supplies no further conditions. The aid's robustness lemma addresses Problem 1's local issue, not a classification for either part of Problem 2.
\subsection{Short English statement}
The main question asks whether generic automorphisms of superstable metric theories are supersimple after suitable perturbations. The supplied aid proves only a finite-orbit-depth continuity bound preserving positive inconsistency gaps; global indiscernibility and orbit coherence remain unresolved. The appended Ehrenfeucht and compact-model-quotient questions are kept as separate problems.
\subsection{Sources}
\begin{itemize}
\item Local UnsolvedMath registry, permanent ID 20002326, source code catalogue entry; canonical \texttt{aim-logic-notes.json}, notes[101].
\item AIM \emph{Model theory of metric structures}, Problems 1 and appended 2: \url{https://aimath.org/WWN/continuouslogic/continuouslogic.pdf}.
\item Supplied background attributions: Pillay--Chatzidakis; Ben-Yaacov; Ben-Yaacov and Berenstein. No full bibliographic entries or exact perturbation definitions are supplied.
\item Registry research attempt 1, dated 2026-08-14: partial result, status at run open, rigor ``proved'', novelty confidence low. The explicitly machine-generated, non-peer-reviewed aid requires independent verification before citation.
\end{itemize}
\end{document}
